% 数学ノート集・共通テンプレート v2（積分論2026年度第1部の版面を参考）
% A5判。章内で定義・定理・例・注意等の番号を共通化。
% コンパイル： platex -interaction=nonstopmode math_note_template_v2.tex (2回)
%              dvipdfmx math_note_template_v2.dvi
\documentclass[a5paper,10pt,onecolumn]{jsarticle}
\usepackage[dvipdfmx,top=22mm,bottom=20mm,inner=19mm,outer=18mm,headheight=13pt,headsep=6mm]{geometry}
\usepackage{amsmath,amssymb,amsthm,mathtools}
\usepackage{enumitem,fancyhdr,titlesec}
\usepackage[dvipdfmx,bookmarks=true,colorlinks=false]{hyperref}
\usepackage{pxjahyper}

% ----- 差し替える文書情報 -----
\newcommand{\NoteTitle}{被覆空間と基本群}
\newcommand{\NoteSeries}{数学ノート集}

% ----- 文字・組版 -----
% 参考PDFの質素な明朝＋Computer Modernの印象を引き継ぎながら，
% A5判で窮屈にならないよう本文の行間を少し確保する．
\linespread{1.13}
\setlength{\parindent}{1zw}
\setlength{\parskip}{0pt}
\setlength{\abovedisplayskip}{6pt plus 2pt minus 2pt}
\setlength{\belowdisplayskip}{6pt plus 2pt minus 2pt}
\setlength{\abovedisplayshortskip}{3pt plus 2pt minus 1pt}
\setlength{\belowdisplayshortskip}{4pt plus 2pt minus 1pt}
\setlist[enumerate]{topsep=3pt,itemsep=1pt,parsep=0pt,leftmargin=2.4zw}
\titleformat{\section}{\normalfont\normalsize\bfseries}{\thesection}{0.8zw}{}
\titlespacing*{\section}{0pt}{2.0ex plus .4ex}{.9ex}
\titleformat{\subsection}{\normalfont\normalsize\bfseries}{\thesubsection}{0.7zw}{}
\titlespacing*{\subsection}{0pt}{1.5ex plus .3ex}{.6ex}
\raggedbottom

% ----- 柱とページ番号：ページ上部、罫線なし -----
\fancypagestyle{mathnote}{%
  \fancyhf{}%
  \fancyhead[L]{\footnotesize \NoteSeries\ （\NoteTitle）}%
  \fancyhead[R]{\footnotesize \thepage}%
  \renewcommand{\headrulewidth}{0pt}%
  \renewcommand{\footrulewidth}{0pt}%
}
\pagestyle{mathnote}

% ----- 節内通し番号の定義・定理群 -----
\newtheoremstyle{mathnotetheorem}
 {4.5pt}{4.5pt}{\normalfont}{}%
 {\bfseries}{}{0.65zw}%
 {\thmname{#1}\thmnumber{ #2}\thmnote{（#3）}}
\theoremstyle{mathnotetheorem}
\newtheorem{theorem}{定理}[section]
\newtheorem{definition}[theorem]{定義}
\newtheorem{lemma}[theorem]{補題}
\newtheorem{proposition}[theorem]{命題}
\newtheorem{corollary}[theorem]{系}
\newtheorem{example}[theorem]{例}
\newtheorem{remark}[theorem]{注意}
\renewcommand{\proofname}{証明}
\renewcommand{\qedsymbol}{$\square$}
\makeatletter
\renewenvironment{proof}[1][\proofname]{%
  \par\pushQED{\qed}\normalfont\topsep4\p@\@plus2\p@\relax
  \trivlist\item[\hskip\labelsep\bfseries #1]\ignorespaces
}{\popQED\endtrivlist\@endpefalse}
\makeatother

% ----- 使うノートで適宜拡張する記号 -----
\newcommand{\R}{\mathbb{R}}
\newcommand{\norm}[1]{\lVert #1\rVert}
\newcommand{\abs}[1]{\lvert #1\rvert}

\begin{document}
\section*{読了後にできること}
本稿を読むための前提は，被覆空間，持ち上げ，基本群．
本稿の具体的な到達課題は，円周の被覆とデッキ変換群を明示する．

\section{被覆写像}

本稿では位相空間はHausdorffであるとし，
特に断らない限り底空間 $X$ は道連結かつ局所道連結とする．

\begin{definition}
連続な全射 $p:E\to X$ が\emph{被覆写像}であるとは，
任意の $x\in X$ に対して開近傍 $U$ が存在し，
\[
 p^{-1}(U)=\bigsqcup_{\lambda\in\Lambda}V_\lambda
\]
と開集合の互いに交わらない和に分解でき，
各制限 $p|_{V_\lambda}:V_\lambda\to U$ が
同相写像になることをいう．
$U$ を\emph{均等被覆近傍}，
$V_\lambda$ をその\emph{シート}という．
\end{definition}

\begin{example}[円周の普遍被覆]
$p:\mathbb R\to S^1$ を $p(t)=e^{2\pi it}$ と定める．
これは被覆写像であり，
$p^{-1}(1)=\mathbb Z$ である．
\end{example}
\begin{proof}
$z=e^{2\pi it_0}\in S^1$ を取る．
$0<\delta<1/2$ として
$U=\{e^{2\pi it}:t_0-\delta<t<t_0+\delta\}$
と置く．
\[
 p^{-1}(U)=
 \bigsqcup_{n\in\mathbb Z}
 (t_0+n-\delta,t_0+n+\delta)
\]
であり，
各区間への $p$ の制限は $U$ への同相写像となる．
最後の主張は $e^{2\pi it}=1$ と $t\in\mathbb Z$ が
同値であることによる．
\end{proof}

\begin{proposition}\label{prop:local-homeo}
被覆写像 $p:E\to X$ は局所同相写像である．
従って $p$ は開写像であり，
各ファイバー $p^{-1}(x)$ は離散部分空間である．
\end{proposition}
\begin{proof}
任意の $e\in E$ について $p(e)$ の均等被覆近傍 $U$ を取り，
$e$ を含むシート $V$ を選ぶと，
$p|_V:V\to U$ は同相写像である．
従って局所同相写像である．
任意の開集合 $W\subset E$ は局所同相写像の
開近傍の和に分解でき，その像は開だから $p$ は開写像である．
各 $e\in p^{-1}(x)$ はシート $V$ とファイバーの交わりが
$\{e\}$ となるので，ファイバーは離散である．
\end{proof}

\section{道の持ち上げ}

\begin{theorem}[道の持ち上げ]\label{thm:path-lift}
$p:E\to X$ を被覆写像，
$\gamma:[0,1]\to X$ を連続な道とする．
$e_0\in p^{-1}(\gamma(0))$ に対し，
ただ一つの連続な道
$\widetilde\gamma:[0,1]\to E$ が存在して
\[
 \widetilde\gamma(0)=e_0,\qquad
 p\circ\widetilde\gamma=\gamma
\]
を満たす．
\end{theorem}
\begin{proof}
$\gamma([0,1])$ はコンパクトなので，
その各点を含む均等被覆近傍の有限部分被覆を取れる．
区間 $[0,1]$ のLebesgue数の補題により，
分割
$0=t_0<\cdots<t_N=1$ を十分細かく取れば，
各 $\gamma([t_{j-1},t_j])$ は
ある均等被覆近傍 $U_j$ に含まれる．
$t_0$ で $e_0$ が与えられている．
$e_0$ を含む $p^{-1}(U_1)$ のシートを $V_1$ とすれば，
\[
 \widetilde\gamma(t)
 =(p|_{V_1})^{-1}(\gamma(t))
 \qquad(t\in[t_0,t_1])
\]
と定められる．
次に $\widetilde\gamma(t_1)$ を含む
$p^{-1}(U_2)$ のシート $V_2$ を取り，
同様に逆写像を用いて
$[t_1,t_2]$ へ延長する．
この操作を有限回繰り返せば，
分割点で値が一致する連続な持ち上げが得られる．

一意性を示す．二つの持ち上げ
$\widetilde\gamma_1,\widetilde\gamma_2$ が
同じ始点を持つとする．
一致する時刻の集合 $A\subset[0,1]$ は，
$E$ がHausdorffなので閉である．
$t\in A$ なら，共通の点
$\widetilde\gamma_1(t)=\widetilde\gamma_2(t)$ を含む
被覆のシート $V$ を取れる．
連続性から $t$ の近傍では両持ち上げが $V$ に入り，
$p|_V$ の単射性から一致する．
従って $A$ は開でもある．
$0\in A$ かつ $[0,1]$ は連結なので
$A=[0,1]$ となる．
\end{proof}

\begin{corollary}\label{cor:closed-lift}
道 $\gamma$ が閉曲線であり，その持ち上げが
$\widetilde\gamma(0)=\widetilde\gamma(1)$ を満たすとき，
$\widetilde\gamma$ は $E$ の閉曲線である．
ただし一般の閉曲線の持ち上げは閉曲線とは限らない．
\end{corollary}
\begin{proof}
前半は閉曲線の定義による．
後半は $p:\mathbb R\to S^1$ と
$\gamma(t)=e^{2\pi it}$ を考えると，
$0$ から始まる持ち上げは
$\widetilde\gamma(t)=t$ となり，
終点は $1\ne0$ であることから分かる．
\end{proof}

\section{ホモトピーの持ち上げ}

\begin{theorem}[ホモトピー持ち上げ定理]\label{thm:homotopy-lift}
$p:E\to X$ を被覆写像，
$H:[0,1]\times[0,1]\to X$ を連続写像とする．
下辺の道 $H(0,t)$ の持ち上げ
$\widetilde h_0:[0,1]\to E$ が与えられたとき，
ただ一つの連続写像
$\widetilde H:[0,1]^2\to E$ が存在して
\[
 p\circ\widetilde H=H,\qquad
 \widetilde H(0,t)=\widetilde h_0(t)
\]
となる．
\end{theorem}
\begin{proof}
まず $[0,1]^2$ の各点について
$H$ の像が入る均等被覆近傍を取る．
$H$ はコンパクトな定義域上で一様連続なので，
$[0,1]^2$ を十分細かい格子へ分割し，
各閉小正方形の像が一つの均等被覆近傍に含まれるようにする．
以下，格子の下段から順に持ち上げを構成する．

最下段の下辺では $\widetilde h_0$ が与えられている．
各小正方形 $Q$ の下辺または左辺の持ち上げが既に得られていて，
その一つの点 $e$ が指定されているとする．
$H(Q)\subset U$ となる均等被覆近傍 $U$ を取り，
$e$ を含むシート $V$ を選ぶ．
\[
 \widetilde H|_Q=(p|_V)^{-1}\circ H|_Q
\]
と置くと，既に定めた辺の持ち上げと
道の持ち上げの一意性によって一致する．
格子を左から右，次に下から上へ進めれば，
有限回で全ての小正方形上に持ち上げを定められる．
隣接する小正方形の共通辺では一意性によって値が一致するので，
貼り合わせの補題から全体の持ち上げは連続である．

一意性については，任意の $t$ に対し
$s\mapsto H(s,t)$ の持ち上げは
$s=0$ における値 $\widetilde h_0(t)$ で一意に決まる．
従って二つの持ち上げは全ての $(s,t)$ で一致する．
\end{proof}

\begin{corollary}[端点固定ホモトピーの持ち上げ]\label{cor:homotopy-end}
始点と終点を固定した二つの道
$\gamma_0,\gamma_1$ が端点固定ホモトピーで結ばれ，
それらの持ち上げが同じ点から始まるなら，
二つの持ち上げの終点は一致する．
\end{corollary}
\begin{proof}
定理\ref{thm:homotopy-lift}によって
ホモトピー $H$ を持ち上げる．
$H(s,0)$ は一定なので，その持ち上げ
$\widetilde H(s,0)$ も道の持ち上げの一意性から一定である．
同様に $H(s,1)$ は一定だから
$\widetilde H(s,1)$ は一定である．
特に $s=0,1$ における終点が一致する．
\end{proof}

\section{被覆と基本群}

\begin{theorem}[被覆が誘導する基本群の単射]
\label{thm:pi1-inj}
$p:(E,e_0)\to(X,x_0)$ を被覆写像とする．
このとき誘導準同型
\[
 p_*:\pi_1(E,e_0)\to\pi_1(X,x_0)
\]
は単射である．
\end{theorem}
\begin{proof}
$[\alpha]\in\pi_1(E,e_0)$ が $p_*[\alpha]=1$ を満たすとする．
$\gamma=p\circ\alpha$ は $X$ で定数曲線へ端点固定ホモトピーで
縮められる．
そのホモトピーを $\alpha$ の初期持ち上げから
定理\ref{thm:homotopy-lift}によって持ち上げる．
端点を固定したホモトピーの持ち上げでは
左右端の像も一定であり，
最後の曲線は $e_0$ における定数曲線の持ち上げとなる．
よって $\alpha$ は $E$ 内で定数曲線へ
端点固定ホモトピーで縮められる．
従って $[\alpha]=1$ であり，$p_*$ は単射である．
\end{proof}

\begin{proposition}[閉曲線への持ち上げ条件]\label{prop:loop-lift}
$p:(E,e_0)\to(X,x_0)$ を被覆写像とし，
$E$ は道連結であるとする．
基点 $x_0$ の閉曲線 $\gamma$ を
$e_0$ から持ち上げた曲線が閉じることと，
\[
 [\gamma]\in p_*(\pi_1(E,e_0))
\]
とは同値である．
\end{proposition}
\begin{proof}
$\gamma$ の持ち上げ $\widetilde\gamma$ が閉じるなら
$[\gamma]=p_*[\widetilde\gamma]$ である．
逆に $[\gamma]=p_*[\alpha]$ となる
$E$ の基点付き閉曲線 $\alpha$ があれば，
$\gamma$ と $p\circ\alpha$ は端点固定ホモトピーで結ばれる．
$e_0$ からの持ち上げの終点は
系\ref{cor:homotopy-end}により等しい．
$p\circ\alpha$ の持ち上げは $\alpha$ なので，
その終点は $e_0$ であり，$\gamma$ の持ち上げも閉じる．
\end{proof}

\section{普遍被覆の構成}

\begin{definition}
$X$ が\emph{半局所単連結}であるとは，
任意の $x\in X$ に対して開近傍 $U$ が存在し，
包含 $U\hookrightarrow X$ が誘導する
$\pi_1(U,x)\to\pi_1(X,x)$ が零準同型となることをいう．
\end{definition}

以下，$X$ は道連結・局所道連結・半局所単連結とする．
局所道連結性によって，各点は
$X$ で零ホモトピーとなる全ての閉曲線しか持たない
道連結な開近傍を持つ．
そのような開集合を\emph{小さい開集合}と呼ぶ．

\begin{definition}
基点 $x_0\in X$ を固定する．
$x_0$ から始まる道 $\gamma$ の端点固定ホモトピー類を
$[\gamma]$ と書き，
\[
 \widetilde X=
 \{[\gamma]:\gamma(0)=x_0\}
\]
と置く．
$p:\widetilde X\to X$ を
$p([\gamma])=\gamma(1)$ と定める．
\end{definition}

\begin{lemma}[普遍被覆の局所座標]\label{lem:universal-local}
$U\subset X$ を小さい開集合，
$[\gamma]\in\widetilde X$ を終点 $x\in U$ の元とする．
\[
 B([\gamma],U)=
 \{[\gamma*\eta]:
  \eta\text{ は }U\text{ 内の }x\text{ からの道}\}
\]
と置く．このとき
\[
 p|_{B([\gamma],U)}:B([\gamma],U)\to U
\]
は全単射である．
また $p^{-1}(U)$ はこの形の集合の互いに交わらない和に
分解される．
\end{lemma}
\begin{proof}
$U$ は道連結なので，任意の $y\in U$ に対して
$x$ から $y$ への道 $\eta$ が存在する．
従って $p|_B$ は全射である．
$\eta_1,\eta_2$ がともに $x$ から $y$ への
$U$ 内の道なら
$\eta_1*\eta_2^{-1}$ は $U$ 内の閉曲線である．
$U$ が小さいので，この閉曲線は $X$ 内で
零ホモトピーである．
従って $[\gamma*\eta_1]=[\gamma*\eta_2]$ となり，
$p|_B$ は単射である．
同じ $U$ について
$B([\gamma],U)$ と $B([\delta],U)$ が一点で交われば，
両者の基準となる経路の差が $U$ 内の道を介して
$X$ 内で零ホモトピーとなるため二つの集合は一致する．
交わらなければ互いに素である．
任意の $[\eta]\in p^{-1}(U)$ は
$B([\eta],U)$ に属するので，
$p^{-1}(U)$ はこのような集合の互いに交わらない和となる．
\end{proof}

\begin{theorem}[普遍被覆の存在]\label{thm:universal-cover}
$X$ を道連結・局所道連結・半局所単連結なHausdorff空間とする．
このとき $X$ は道連結かつ単連結な被覆空間
$p:\widetilde X\to X$ を持つ．
\end{theorem}
\begin{proof}
補題\ref{lem:universal-local}の集合
$B([\gamma],U)$ を開集合の基として
$\widetilde X$ に位相を定める．
この集合族が基になることを確認する．
任意の $[\gamma]\in\widetilde X$ は
終点の小さい開近傍 $U$ に対して
$B([\gamma],U)$ に属する．
二つの基本集合の共通部分に属する
$[\delta]$ を取れば，その終点に含まれる
十分小さい道連結な小さい開集合
$W\subset U\cap V$ が存在する．
$B([\delta],W)$ は二つの基本集合の共通部分に含まれる．
従って位相の基の公理が成立する．
また $p|_{B([\gamma],U)}$ は全単射であり，
基本集合の定義からその逆写像は
$U$ の十分小さい道連結開集合上で連続である．
実際，$W\subset U$ が道連結で小さければ
$p^{-1}(W)\cap B([\gamma],U)$ は
同じ形式の基本集合になる．
従って $p|_B$ は同相写像である．
補題\ref{lem:universal-local}から
$p^{-1}(U)$ は互いに交わらないシートの和となるので
$p$ は被覆写像である．

道連結性を示す．任意の $[\gamma]\in\widetilde X$ に対し
\[
 \widetilde\gamma(t)=
 [\gamma|_{[0,t]}]
\]
と置く．
局所座標の定義から $\widetilde\gamma$ は連続であり，
空の道の類 $[\mathrm{const}_{x_0}]$ と
$[\gamma]$ を結ぶ．
従って $\widetilde X$ は道連結である．

最後に単連結性を示す．
$\widetilde x_0=[\mathrm{const}_{x_0}]$ とする．
$\widetilde X$ 内の基点付き閉曲線 $\alpha$ を取り，
$\gamma=p\circ\alpha$ と置く．
$\alpha$ は $\gamma$ の $\widetilde x_0$ からの
持ち上げなので，経路の持ち上げの一意性から
\[
 \alpha(t)=[\gamma|_{[0,t]}]
\]
である．
$\alpha(1)=\widetilde x_0$ だから
$[\gamma]=[\mathrm{const}_{x_0}]$ となり，
$\gamma$ は $X$ 内で零ホモトピーである．
定理\ref{thm:pi1-inj}により
$p_*:\pi_1(\widetilde X,\widetilde x_0)\to\pi_1(X,x_0)$
は単射であり，$p_*[\alpha]=[\gamma]=1$ だから
$[\alpha]=1$ である．
ゆえに $\widetilde X$ は単連結である．
\end{proof}

\section{被覆空間の分類}

以下，$G=\pi_1(X,x_0)$ とする．

\begin{definition}
$H\le G$ を部分群とする．
$x_0$ から始まる二つの道 $\gamma,\delta$ に対し
\[
 \gamma\sim_H\delta
 \quad\Longleftrightarrow\quad
 \gamma(1)=\delta(1),\
 [\gamma*\delta^{-1}]\in H
\]
と定める．
この同値類を $[\gamma]_H$ と書き，
$X_H$ をその全体とする．
$p_H:X_H\to X$ を
$p_H([\gamma]_H)=\gamma(1)$ と定める．
\end{definition}

\begin{lemma}\label{lem:h-cover}
$\sim_H$ は同値関係である．
$X_H$ に，普遍被覆と同じく小さい開集合 $U$ に対する
\[
 B_H([\gamma]_H,U)
 =\{[\gamma*\eta]_H:
 \eta\text{ は }\gamma(1)\text{ からの }U\text{ 内の道}\}
\]
を基とする位相を入れると，
$p_H:X_H\to X$ は道連結な被覆写像になる．
\end{lemma}
\begin{proof}
反射律は定数曲線の類が $H$ に属すること，
対称律は群の逆元に対する閉性，
推移律は閉曲線類の積に対する閉性から従う．
具体的に
$[\gamma*\delta^{-1}]\in H$，
$[\delta*\eta^{-1}]\in H$ なら，
その積は $[\gamma*\eta^{-1}]$ に等しいので推移律を得る．

$U$ 内の二つの道が同じ始点と終点を持つなら，
その差は $X$ 内で零ホモトピーとなるため
$\sim_H$ に関して同値である．
従って $p_H|_{B_H([\gamma]_H,U)}$ は全単射である．
基本集合の交わりと局所的な逆写像の連続性は
補題\ref{lem:universal-local}および
定理\ref{thm:universal-cover}の証明と同様の
小さい開集合の議論で従う．
よって $p_H$ は被覆写像である．
また各 $[\gamma]_H$ は始点を空の道の類とする
部分道の列 $t\mapsto[\gamma|_{[0,t]}]_H$ によって
結ばれるため $X_H$ は道連結である．
\end{proof}

\begin{theorem}[基点付き被覆の分類]\label{thm:cover-class}
$X$ を道連結・局所道連結・半局所単連結なHausdorff空間とする．
基点付きの道連結被覆
$p:(E,e_0)\to(X,x_0)$ の同型類と，
$\pi_1(X,x_0)$ の部分群全体の間には自然な全単射がある．
対応は
\[
 (E,e_0)\longmapsto
 p_*(\pi_1(E,e_0))
\]
で与えられる．
\end{theorem}
\begin{proof}
まず $H\le G$ から補題\ref{lem:h-cover}で
被覆 $p_H:X_H\to X$ を構成する．
基点を $[\mathrm{const}_{x_0}]_H$ とする．
基点付き閉曲線 $\gamma$ の持ち上げが閉じることは
$[\gamma]_H=[\mathrm{const}_{x_0}]_H$ と同値であり，
これは $[\gamma]\in H$ と同値である．
命題\ref{prop:loop-lift}から
$(p_H)_*\pi_1(X_H)=H$ となる．

逆に基点付き道連結被覆
$p:(E,e_0)\to(X,x_0)$ を取る．
$H=p_*\pi_1(E,e_0)$ と置く．
各 $[\gamma]_H\in X_H$ に対し，
$\gamma$ の $e_0$ からの唯一の持ち上げを
$\widetilde\gamma$ とし
\[
 \Phi([\gamma]_H)=\widetilde\gamma(1)
\]
と定める．
これがwell-definedであることを示す．
$\gamma\sim_H\delta$ なら
$\gamma*\delta^{-1}$ の基点 $e_0$ からの持ち上げは
命題\ref{prop:loop-lift}によって閉じる．
従って $\gamma$ と $\delta$ の持ち上げの終点は一致する．
逆に持ち上げの終点が一致するなら，
$\gamma*\delta^{-1}$ の持ち上げは閉じるので
$[\gamma*\delta^{-1}]\in H$ となる．
従って $\Phi$ は単射である．
$E$ は道連結なので，任意の $e\in E$ は
$e_0$ からの道 $\widetilde\gamma$ の終点であり，
$\Phi([\gamma]_H)=e$ となるから全射である．

最後に $\Phi$ が同相であることを示す．
任意の $[\gamma]_H$ の終点 $x$ に，
$E$ について均等被覆され，かつ
$X$ の小さい道連結開集合である近傍 $U$ を取る．
$\Phi$ は $B_H([\gamma]_H,U)$ を，
$\widetilde\gamma(1)$ を含む $p^{-1}(U)$ のシートへ写す．
この対応は両方の空間で $U$ への被覆射影と可換であり，
両射影の制限は同相なので $\Phi$ も局所的に同相となる．
$\Phi$ は全単射なので全体でも同相である．
また $p\circ\Phi=p_H$ を満たし，基点を保つ．
従って被覆は $X_H$ と基点付き同型である．
以上から全単射を得る．
\end{proof}

\begin{corollary}[基点を忘れた分類]\label{cor:unbased}
基点を指定しない道連結被覆の同型類は，
$\pi_1(X,x_0)$ の部分群の共役類に対応する．
\end{corollary}
\begin{proof}
同じ被覆 $p:E\to X$ のファイバー $p^{-1}(x_0)$ に属する
二点 $e_0,e_1$ を取る．
$E$ は道連結なので $e_0$ から $e_1$ への道 $\alpha$ がある．
$p\circ\alpha$ は $x_0$ における閉曲線であり，
基点変更の標準同型によって
\[
 p_*(\pi_1(E,e_1))
 =[p\circ\alpha]^{-1}
 p_*(\pi_1(E,e_0))
 [p\circ\alpha]
\]
となる．
従って基点を変えると部分群は共役に変わる．
逆に任意の $g\in\pi_1(X,x_0)$ を表す閉曲線を
$e_0$ から持ち上げ，その終点を $e_1$ とすれば，
対応する部分群は $g$ による共役となる．
基点付き分類定理を適用すれば結論が従う．
\end{proof}

\section{被覆変換群}

\begin{definition}
被覆 $p:E\to X$ の\emph{被覆変換}とは，
同相写像 $\varphi:E\to E$ で
$p\circ\varphi=p$ を満たすものをいう．
全ての被覆変換は合成に関して群をなし，
$\operatorname{Deck}(E/X)$ と書く．
\end{definition}

\begin{lemma}[被覆変換の一意性]\label{lem:deck-unique}
$E$ が道連結であるとき，
二つの被覆変換が一点で一致すれば全体で一致する．
従って被覆変換のファイバーへの作用は自由である．
\end{lemma}
\begin{proof}
$\varphi,\psi$ が $e_0$ で一致するとする．
任意の $e\in E$ を取り，
$e_0$ から $e$ への道 $\gamma$ を選ぶ．
$p\circ\varphi\circ\gamma
=p\circ\gamma
=p\circ\psi\circ\gamma$
であり，二つの道は始点で一致する．
道の持ち上げの一意性によって
$\varphi\circ\gamma=\psi\circ\gamma$ となる．
終点で $\varphi(e)=\psi(e)$ が従う．
\end{proof}

\begin{theorem}[被覆変換群と正規化群]\label{thm:deck-normalizer}
$H=p_*\pi_1(E,e_0)\le G=\pi_1(X,x_0)$ とする．
このとき
\[
 \operatorname{Deck}(E/X)\cong N_G(H)/H
\]
が成り立つ．ただし $N_G(H)=\{g\in G:gHg^{-1}=H\}$ とする．
特に $H$ が正規部分群なら被覆変換群は $G/H$ と同型である．
\end{theorem}
\begin{proof}
分類定理によって $E$ を $X_H$ と同一視する．
$g=[\alpha]\in G$ に対し，
\[
 \Phi_g([\gamma]_H)=[\alpha*\gamma]_H
\]
と定めることを考える．
$\gamma\sim_H\delta$ なら
$[\gamma*\delta^{-1}]\in H$ である．
従って $\alpha*\gamma$ と $\alpha*\delta$ が
$\sim_H$ に関して同値である条件は
$g[\gamma*\delta^{-1}]g^{-1}\in H$ である．
よって $\Phi_g$ がwell-definedとなり，
かつ逆写像 $\Phi_{g^{-1}}$ もwell-definedとなるための
必要十分条件は $gHg^{-1}=H$，すなわち $g\in N_G(H)$ である．
このとき $\Phi_g$ は被覆変換である．
また $\Phi_g$ が恒等写像となることは
$\Phi_g([\mathrm{const}]_H)=[\mathrm{const}]_H$ と同値であり，
$g\in H$ と同値である．

逆に任意の被覆変換 $\Phi$ を取る．
$\Phi([\mathrm{const}]_H)$ は $x_0$ 上のファイバーの元なので，
ある基点付き閉曲線 $\alpha$ に対して
$[\alpha]_H$ と書ける．
補題\ref{lem:deck-unique}から
$\Phi$ はファイバーの一点での値によって一意に決まる．
この $\alpha$ による写像 $\Phi_g$ が
被覆変換として定まるためには，
被覆変換が基本群の部分群を保つことから
$gHg^{-1}=H$ が必要である．
従って $\Phi=\Phi_g$ であり，
準同型 $N_G(H)\to\operatorname{Deck}(E/X)$ は全射で核が $H$ となる．
第一同型定理から結論を得る．
\end{proof}

\section{基本的な例}

\begin{example}[円周の連結被覆の分類]\label{ex:circle}
$\pi_1(S^1,1)\cong\mathbb Z$ なので，
$S^1$ の道連結な被覆は，
各 $n\ge1$ に対する $z\mapsto z^n:S^1\to S^1$ と
普遍被覆 $\mathbb R\to S^1$ に分類される．
\end{example}
\begin{proof}
$\mathbb Z$ の部分群は $n\mathbb Z$（$n\ge0$）に限る．
$n\ge1$ のとき $z\mapsto z^n$ の誘導準同型は
$\pi_1(S^1)\cong\mathbb Z$ 上で整数の $n$ 倍となるので
像は $n\mathbb Z$ である．
$n=0$ に対応する零部分群を持つ被覆は
$\mathbb R\to S^1$ である．
分類定理によってこれら以外の連結被覆はない．
\end{proof}

\begin{example}[トーラスの被覆]\label{ex:torus}
$T^2=S^1\times S^1$ について
$\pi_1(T^2)\cong\mathbb Z^2$ である．
従って連結被覆は $\mathbb Z^2$ の部分群に対応する．
特に普遍被覆は $\mathbb R^2\to T^2$ である．
\end{example}
\begin{proof}
直積空間の基本群の公式から
$\pi_1(T^2)\cong\pi_1(S^1)\times\pi_1(S^1)
\cong\mathbb Z^2$ である．
$\mathbb R^2\to T^2$，
$(s,t)\mapsto(e^{2\pi is},e^{2\pi it})$ は
二つの円周被覆の直積であり，均等被覆近傍は
各因子の均等被覆近傍の積として取れる．
$\mathbb R^2$ は単連結なのでこれは普遍被覆である．
\end{proof}

\begin{example}[実射影空間]\label{ex:projective}
$n\ge2$ のとき，
\[
 S^n\to\mathbb RP^n,\qquad x\mapsto\{\pm x\}
\]
は二枚被覆であり，
$\pi_1(\mathbb RP^n)\cong\mathbb Z/2\mathbb Z$ である．
\end{example}
\begin{proof}
各 $x\in S^n$ の小さい開球を
$x$ と $-x$ の互いに交わらない近傍として選べば，
それらの射影は同じ開集合へ同相に写るため被覆写像である．
$n\ge2$ では $S^n$ は単連結なので普遍被覆である．
被覆変換群は恒等写像と対蹠写像 $x\mapsto-x$ の
二元からなる．
普遍被覆の場合は定理\ref{thm:deck-normalizer}において
$H=\{1\}$ なので
$\pi_1(\mathbb RP^n)\cong\operatorname{Deck}(S^n/\mathbb RP^n)
\cong\mathbb Z/2\mathbb Z$ となる．
\end{proof}

\begin{remark}
被覆空間の分類は，
局所的には同相である空間の大域的な違いを
基本群の部分群によって記述する定理である．
半局所単連結性を外すと普遍被覆が存在しない場合があり，
一般にはこの分類の形式はそのまま成立しない．
\end{remark}

\section{到達課題の完全解答}
\noindent\textbf{実直線から円周への普遍被覆．}
\begin{proposition}
$p:\mathbb R\to S^1$，
$p(t)=e^{2\pi it}$ は被覆写像であり，
デッキ変換群は
\[
\operatorname{Deck}(p)
=\{t\mapsto t+n:n\in\mathbb Z\}
\cong\mathbb Z
\]
である．
\end{proposition}
\begin{proof}
円周の各点を含む十分短い開弧 $U$ の逆像は
互いに素な開区間の和であり，
各区間上の $p$ は $U$ への同相写像である．
従って被覆写像となる．
$p(h(t))=p(t)$ を満たすデッキ変換 $h$ について
$h(t)-t\in\mathbb Z$ である．
連続関数 $h(t)-t$ は連結空間 $\mathbb R$ から
離散空間 $\mathbb Z$ への写像なので定数である．
逆に各整数だけの平行移動は被覆を保存する同相写像である．
\end{proof}
\noindent\textbf{到達確認．}
$p(t)=e^{2\pi it}$ における経路
$\gamma(s)=e^{2\pi i3s}$ の0からの持ち上げを求めよ．
\par\noindent\textbf{解答．}
$\widetilde\gamma(s)=3s$ であり，
終点の値3が巻き数を表す．
\end{document}